%% env_description.tex
%% Description of environments, observation spaces, performance metrics, and shift axes.
%% Clarifies experimental scope and eliminates unsupported claims.

\paragraph{Experimental Environments and Evaluation Protocol.}
The primary benchmark domain used throughout this work is the \textbf{Visual Pursuit} environment.
In Visual Pursuit, an agent receives a $64\times 64$ pixel RGB observation representing an arena containing a moving visual target.
The agent outputs continuous 2D actions $(dx, dy) \in [-1, 1]^2$ governing the velocity of a tracking reticle.
The task objective is visual pursuit: the agent must continuously track and center its reticle on the dynamically moving target based purely on image inputs.
Performance is quantified by the \textbf{Center-of-target Localization Error (CLE)}, defined as the Euclidean distance (in pixels) between the tracking reticle center and the ground-truth target centroid at each frame.
An episode is classified as successful if the mean CLE over the episode satisfies $\text{CLE} < 5\text{ px}$.
For context, an untrained random policy achieves a baseline performance floor of approximately $\sim 22\text{ px}$ mean CLE.

\paragraph{Distribution Shift Axes.}
To systematically characterize competence degradation, we subject Visual Pursuit to three parameterized visual shift axes:
\begin{enumerate}
    \item \textbf{Photometric Shifts:} Continuous variations in observation brightness and contrast levels.
    \item \textbf{Camera Noise:} Zero-mean additive Gaussian sensor noise scaled across noise variances $\sigma^2 \in [0, 1]$.
    \item \textbf{Severe Shift:} A composite catastrophic perturbation combining full color inversion, substantial contrast compression, and high-variance Gaussian camera noise, designed to stress-test catastrophic failure modes and floor-collapse behavior.
\end{enumerate}

\paragraph{Scope Clarifications and Removed Domains.}
We explicitly clarify two domain constraints:
\begin{itemize}
    \item \textbf{Procgen CoinRun:} While Procgen CoinRun (a discrete-action platformer game) is analyzed specifically in our Phase~1 representation correlation sanity checks to evaluate semantic weighting ($\tilde{d}_{rep}$) under discrete game mechanics, \emph{Procgen was NOT evaluated for continuous competence boundary forecasting}. Any claims suggesting Procgen was used for competence boundary forecasting have been removed.
    \item \textbf{Reacher Environment:} There is \emph{no} Reacher or robotic arm reaching environment in our experimental results. All prior mentions or placeholder references to Reacher tasks have been excised.
\end{itemize}
